\documentclass[11pt,letterpaper]{article}
\usepackage[margin=1in]{geometry}
\usepackage[T1]{fontenc}
\usepackage{lmodern,microtype,amsmath,amssymb,amsthm,mathtools}
\usepackage[hidelinks]{hyperref}
\setlength{\emergencystretch}{2em}
\numberwithin{equation}{section}
\newtheorem{theorem}{Theorem}[section]
\newtheorem{proposition}[theorem]{Proposition}
\newtheorem{lemma}[theorem]{Lemma}
\theoremstyle{definition}
\newtheorem{example}[theorem]{Example}
\DeclareMathOperator{\Scal}{Scal}
\DeclareMathOperator{\vol}{Vol}
\DeclareMathOperator{\supp}{supp}
\title{Marked immersed index-two handles\\and exterior chiral Weyl bounds}
\author{Edward B. Baker III}
\date{25 November 2026}

\begin{document}
\maketitle
\begin{abstract}
We construct a rounded index-two handle in Euclidean five-space whose
actual linear height realizes the product attaching and belt markings.
Flattening its graph makes both induced chiral Weyl energies decay
quartically in the slope parameter. The estimate extends to two
prescribed compatible, scaled graph collars and their corner overlap.
Fixing a nonflat quadratic collar forces a positive cost for every
smooth immersed filling, including those that cease to be graphs.
A relative Morse calculation forces an index-two point in almost every
direction of the reference normal image, and the indexed Gauss area
formula converts this count to a strict chiral bound. For finitely many
such collars on a closed smooth homotopy four-sphere, the overlap
of their directional images yields a recognition criterion entirely
in terms of the common exterior and the marked references. A signed
count gives an additional necessary budget for negative Gauss mass
on this exterior. Three exact marked collars on a standard sphere
realize the quartic order of the seam cost in a positive Yamabe class.
\end{abstract}

\section{Closed and relative height counts}

Every smooth homotopy four-sphere admits a smooth immersion into
\(\mathbb R^5\). Its tangent bundle is stably trivial and stabilization
in rank five is already in the stable range over a four-dimensional
complex. The existence assertion then follows from the immersion
theorem \cite[Theorem 3.1]{KervaireMilnor1963Groups}
\cite[\S6]{Hirsch1959Immersions}. Prescribing a graph on a marked handle
imposes relative geometric conditions on that immersion. Our aim is to
measure which conditions survive arbitrary replacement of its interior,
and to specify the exterior data needed to use an energy estimate on
the original closed smooth domain.

For an oriented smooth hypersurface immersion \(F\) use the induced
Riemannian metric \(g\), the compatible unit normal \(\nu\), and the
shape operator \(d\nu=-dF\circ A\). Let \(U_2\) be the open set where
\(A\) has two positive and two negative eigenvalues. On a measurable
set \(P\) of the smooth domain put
\[
 \begin{split}
 E_\pm(P)&=\int_P |W_g^\pm|_{\mathrm{HS}}^2\,dV_g,\qquad
 I_2(P)=\int_{P\cap U_2}\det A\,dV_g,\\
 J(P)&=\int_P(-\det A)_+\,dV_g,\qquad
 B^2(P)=\int_P|A-\tfrac14(\operatorname{tr}A)I|^4\,dV_g .
 \end{split}
\]
The Weyl blocks are endomorphisms of the real chiral two-form
bundles. Distinct preimages of an immersed point are counted on the
domain. These integrals are invariant under ambient dilation.
Reversing orientation exchanges the Weyl blocks, and the Gauss
equation gives their pointwise squared norms equal in the Euclidean
hypersurface setting.

\begin{lemma}[Middle-signature curvature and height]\label{lem:area}
Put \(\omega_4=\sigma_4(S^4)=8\pi^2/3\). For almost every
\(u\in S^4\), the critical points of \(\langle F,u\rangle\) in the
interior of \(P\) are nondegenerate and lie in \(\nu^{-1}(\{u,-u\})\).
If \(m_{2,P}(u)\) counts those of index two, then for a smooth compact
domain, or a domain cut out by smooth seams,
\[
 I_2(P)=\tfrac12\int_{S^4}m_{2,P}(u)\,d\sigma_4(u),
 \qquad E_\pm(P)\ge\tfrac83 I_2(P).
\]
If \(P^\circ\cap U_2\ne\varnothing\), the energy comparison is strict.
\end{lemma}

\begin{proof}
At a critical point the height Hessian is \(\pm g(A\cdot,\cdot)\),
according to the normal sign, so the index is two precisely on
\(U_2\). The signed normal Jacobian is \(\det A\). Sard's theorem
and the area formula give the count, with a factor of two from the
two signs. Normal images of smooth seams have spherical measure
zero. This is the indexed form of the total curvature calculation
\cite{ChernLashof1957Curvature,OntiVlachos2017Hypersurfaces}.

In a principal orthonormal frame the Gauss curvature eigenvalues
are \(k_i k_j\). Write \(s_1,s_2,s_3\) for the three opposite-pair
sums, and \(\bar s\) for their mean. Either chiral Weyl block has
eigenvalues \((s_j-\bar s)/2\), so its squared norm is
\(\frac1{12}\sum_{j<\ell}(s_j-s_\ell)^2\).
On \(U_2\), label the curvatures \((-a,-b,c,d)\) with positive
\(a,b,c,d\). The sums are \(ab+cd,-ac-bd,-ad-bc\).
By the arithmetic--geometric mean inequality the first exceeds
or equals \(2\sqrt{\det A}\) and the latter two are at most
\(-2\sqrt{\det A}\). This gives \(\frac83\det A\) as a lower
bound and equality requires \(a=b=c=d\).

If equality held on an open set, the two principal distributions
would have rank two and eigenvalues \(k,-k\) with \(k>0\). For
orthonormal vectors \(X,Y\) in the same distribution, Codazzi gives
\(X(\pm k)=\langle(\nabla_X A)Y,Y\rangle
=\langle(\nabla_Y A)X,Y\rangle=0\).
Thus \(dk=0\). For \(X,Y\) in opposite distributions, the two
sides of Codazzi lie in opposite distributions and separately
vanish. Metric compatibility then makes both distributions
parallel. Every mixed sectional curvature would be zero, whereas
Gauss gives \(-k^2\). The nonnegative continuous pointwise defect
is therefore positive on some open subset of every nonempty
open middle-signature region.
\end{proof}

\begin{lemma}[A smooth handle-count bound]\label{lem:top}
Every Morse function on a closed smooth homotopy four-sphere
not diffeomorphic to \(S^4\) has at least two index-two points.
\end{lemma}

\begin{proof}
Reorder the handles by index
\cite[Theorems 4.7--4.8]{Milnor1965HCobordism}.
The union of the 0- and 1-handles
is connected since all subsequent attaching regions are connected.
Joining the 0-handles along a spanning tree of the 1-handles
gives one smooth four-ball, carrying the remaining attachments to
its boundary \cite[Theorem~5.4]{Milnor1965HCobordism}.
The dual operation gives one 4-handle. These reductions
preserve the number \(m_2\) of index-two handles. Euler
characteristic gives \(m_1+m_3=m_2\) for the reduced counts.
If \(m_2=0\), the two balls glue to \(S^4\), by extension of
boundary diffeomorphisms over the ball \cite{Hatcher1983Smale}.
If \(m_2=1\), reverse the handles if needed to obtain a 0-, a
2-, a 3- and a 4-handle. The boundary after the 2-handle is
\(S^1\times S^2\) by the dual 1-handle description. Hence its
integral surgery coefficient is zero and Property~R identifies the
attaching knot as the unknot \cite[Corollary 8.6]{Gabai1987FoliationsIII}.
The \(S^2\times D^2\) trace and the dual \(S^1\times D^3\)
glue to \(S^4\) with the given boundary identification by the
extension theorem for four-dimensional 1-handlebodies
\cite[Theorem A]{LaudenbachPoenaru1972Handlebodies}.
\end{proof}

In particular a closed immersed homotopy sphere with
\(E_+\le\frac83\omega_4=64\pi^2/9\) is smoothly standard.
On a putative nonstandard type, if \(E_+<8\pi^2\), the mean number
of index-two points is strictly less than \(9/4\). Since each
Morse height has at least two, more than three quarters of directions
have exactly two. To turn this count into a surgery argument one
needs the attaching maps of the same height.

\section{A quadratic handle with two marked faces}
\label{sec:graph}

Write \(z=(x,y)\in\mathbb R^2\times\mathbb R^2\),
\[
 q(z)=(|y|^2-|x|^2)/2,\qquad s(z)=(|y|^2+|x|^2)/2,\qquad
 D=\operatorname{diag}(-1,-1,1,1).
\]
The product \(D^2_x\times D^2_y\) has attaching face \(|x|=1\),
outgoing face \(|y|=1\), and common corner \(S^1\times S^1\).
Choose a smooth even convex function \(\rho\) by convolving absolute
value with a nonnegative even mollifier supported in
\((-\delta_0,\delta_0)\). Then \(|q|\le\rho(q)\le |q|+\delta_0\),
\(|\rho'|\le1\), and \(\rho(q)=|q|\) for \(|q|\ge\delta_0\).
The domain \(\Omega=\{s+\rho(q)\le1\}\) rounds the corner.
Near the rounded portion both radii are positive, and \(s\) is
a transverse coordinate. Elsewhere the boundary agrees with
one of the original faces. The domain is star-shaped since
\(s+\rho(q)\) is nondecreasing in \(|x|^2\) and \(|y|^2\).

\begin{proposition}[Actual height and product framing]\label{prop:trace}
For \(R,t>0\) the map
\[
                 F_{R,t}(z)=(Rz,Rtq(z))
\]
embeds \(\Omega\) in \(\mathbb R^5\). Its actual vertical height is
\(Rtq\) and has just one interior critical point, of index two.
For \(\delta_0<\lambda<1/2\) the trace \(|q|\le\lambda\)
has marked incoming and outgoing solid tori with actual heights
\(-Rt\lambda\) and \(Rt\lambda\). They have the product framing
of the index-two handle and its belt.
\end{proposition}

\begin{proof}
The horizontal differential is \(RI\), and the critical point at
zero has height Hessian \(RtD\). At the given levels the parametrizations
in the domain are
\[
 (\theta,v)\longmapsto(\sqrt{2\lambda+|v|^2}\,\theta,v),
 \qquad
 (w,\phi)\longmapsto(w,\sqrt{2\lambda+|w|^2}\,\phi),
\]
where \(\theta,\phi\in S^1\) and \(|v|,|w|\le\sqrt{1-2\lambda}\).
On the lateral boundary between the levels the coordinates are
\((q,\theta,\phi)\). The two radii there are
\(\sqrt{1-\rho(q)-q}\) and \(\sqrt{1-\rho(q)+q}\),
both positive. In particular \(q\) is a submersion on the lateral
boundary, and the product side identifies both angle coordinates
between the two faces.

In the induced metric the negative height gradient is a positive
multiple of \((x,-y)\). The descending disk is \(y=0\) and the
ascending disk is \(x=0\). At the incoming circle \(v=0\) the
coordinate normal vectors \(\partial_{v_j}\) are the original
\(y\)-vectors. At the outgoing belt circle \(w=0\) the \(x\)-vectors
give the belt marking in the same way. This specifies the
framed solid-torus embeddings in the actual height levels.
Their images in a surrounding boundary are determined by the
supplied exterior attachment.
\end{proof}

\subsection{Slope, normal directions and energy}

Use the upward normal and set \(\mathcal W=(1+t^2|z|^2)^{1/2}\).
At \(R=1\) the metric is \(I+t^2Dz\otimes Dz\), its volume
density is \(\mathcal W\), and the second fundamental form is
\(tD/\mathcal W\). In an induced orthonormal frame the shape operator
is \(tD+O(t^3)\), uniformly on the marked domain, and has signature
\((2,2)\) even at the seam. At general \(R\),
\[
 \nu=\frac{(-tDz,1)}{\mathcal W},\qquad
 \det A=\frac{t^4}{R^4\mathcal W^6},\qquad
 dV_g=R^4\mathcal W\,dz,\qquad
 \Scal_g=-\frac{2t^2(2+t^2|z|^2)}{R^2\mathcal W^4}.
\]
The scalar identity follows from
\(\Scal=(\operatorname{tr}A)^2-|A|^2\) and the rank-one
inverse formula. In the chiral eigenvalue formula the leading
eigenvalues are \( (4/3,-2/3,-2/3)t^2/R^2\).
For every measurable \(V\subset\Omega\), with \(|V|\) its
horizontal coordinate volume, this gives
\begin{equation}\label{eq:scale}
 \begin{aligned}
 E_\pm(V)&=\tfrac83 |V|t^4+O(|V|t^6),&
 I_2(V)&=|V|t^4+O(|V|t^6),\\
 B^2(V)&=16|V|t^4+O(|V|t^6),& J(V)&=0 .
 \end{aligned}
\end{equation}
The constants are independent of \(R\) and \(V\). In particular
ambient dilation alone leaves the energy fixed.

For a graph of slope at most \(S\), its upward normal directions
lie in the spherical cap about \(e_5\) of angular radius
\(\arctan S\). Its possible critical directions are in that
cap and its antipode. Each of these caps has area
\(|S^3|\int_0^{\arctan S}\sin^3\alpha\,d\alpha=O(S^4)\)
as \(S\downarrow0\). For the reference graph define the precise
one-sign normal image \(N_t=\nu(\Omega^\circ)\). The upper
gnomonic coordinates \(p=u_{\rm hor}/u_5\) give \(p=-tDz\).
The spherical area density in these coordinates is
\((1+|p|^2)^{-5/2}dp\), so
\begin{equation}\label{eq:normal-area}
       s_t:=\sigma_4(N_t)
       =t^4\int_\Omega(1+t^2|z|^2)^{-5/2}\,dz
       =I_2(\Omega)>0 .
\end{equation}
The graph has one index-two point for each direction in
\(N_t\cup(-N_t)\), and zero for directions outside these two
sets and the seam normal images.

\subsection{Compatible graph collars}

Fix a normalized collar width \(\epsilon>0\). Suppose marked
graph germs \((Rz,Rt k_{x,t}(z))\) and \((Rz,Rt k_{y,t}(z))\)
are given on \(|x|>1-\epsilon\) and \(|y|>1-\epsilon\)
in the same horizontal coordinates. Require equality as functions
on the overlap and uniform normalized \(C^2\) bounds. For gluing
to an exterior, use the germs supplied across the same marked seam.
Let \(k_t\) be their union. Choose radial cutoffs \(\chi_x,\chi_y\)
supported in the respective collars, equal to one beyond
radius \(1-\epsilon/2\), with \(j\)-th derivatives
bounded by \(C\epsilon^{-j}\) for \(j\le2\).
Define
\[
 \chi=1-(1-\chi_x)(1-\chi_y),\qquad
 f_t=q+\eta_t,\qquad \eta_t=\chi(k_t-q).
\]
The added function extends smoothly by zero. With
\(\sqrt{1-\delta_0}>1-\epsilon/2\), a whole neighborhood
of the rounded seam has \(\chi=1\). Thus the graph of \(Rt f_t\)
with horizontal map \(Rz\) matches the prescribed collars,
including their common corner, and is an embedding of the patch.

Let \(M_j=\sup |D^j(k_t-q)|\) on the collar union for
\(0\le j\le2\) and set
\[
        K=C(1+M_2+\epsilon^{-1}M_1+\epsilon^{-2}M_0).
\]
Then \(K\) bounds the coordinate Hessian of \(f_t\) there.
At \(R=1\), for a graph of height \(l\),
\(|A|\le|D^2l|/(1+|Dl|^2)^{1/2}\) and
\(dV=(1+|Dl|^2)^{1/2}dz\). The Gauss equation bounds each
of \(|W^\pm|^2,|\det A|,|A^\circ|^4,\Scal_-^2\) by
\(C|A|^4\). On either collar minus the other and on their
overlap, the integral of any of these quantities over \(V\)
is at most \(C|V|t^4K^4\), also at general \(R\) by scale
invariance. These three sets and the core form an additive
partition up to null boundaries. The normalized bounds also give
slope \(O(t)\), so their normal images lie in the same type of
spherical caps of area \(O(t^4)\). The energy estimate is quartic at
fixed widths and records the derivative loss if those widths shrink.

For the actual framing and critical point after matching,
assume in addition \(\|\eta_t\|_{C^1}\) sufficiently small.
Choose \(\sigma>0\) with \(\lambda+2\sigma<1/2\), and
\(1-\epsilon>\sqrt{2(\lambda+2\sigma)}\).
The correction then vanishes near the core disks over
\(|q|\le\lambda+2\sigma\). On its support \(|Dq|\ge1-\epsilon\).
The functions \(q+v\eta_t\), \(0\le v\le1\), preserve the
interior critical point and the boundary submersion near the
two specified levels, since the \(q\)-derivative along the
side of the unperturbed trace is one. On the regular part
of a slightly larger band choose smooth \(Z_v\) with
\(d(q+v\eta_t)(Z_v)=1\), tangent to the side there.
Local submersions give such fields and a partition of unity
joins them. The flow of \(-\eta_t Z_v\), cut off outside
that larger band, fixes the core disk neighborhoods and
transports the two marked levels. The time derivative of
\((q+v\eta_t)\) along this flow vanishes where the cutoff is one.
The supplied side identifications must use this transport
if the side height changes. Exact equality with \(q\) there
keeps the original side parametrization. A bound on the
normalized \(C^2\) jets alone gives the energy estimate,
while the additional submersion control gives this height result.

\section{A positive cost for a fixed collar}\label{sec:fixed}

\begin{theorem}[Fixed quadratic germ]\label{thm:fixed}
Let \(R>0\) and \(t>0\) be fixed. Suppose a smooth immersion
\(G:\Omega\looparrowright\mathbb R^5\) agrees as a map with
\(F_{R,t}\) on a neighborhood of the entire marked rounded boundary.
For almost every \(u\in N_t\cup(-N_t)\), its actual height
\(\langle G,u\rangle\) has an interior index-two critical point.
In particular
\[
                   I_2(G;\Omega)\ge s_t,\qquad
                   E_\pm(G;\Omega)>\tfrac83 s_t .
\]
The estimate allows self-intersections and arbitrary immersed
interiors. Its lower bound is independent of \(R\).
\end{theorem}

\begin{proof}
Fix \(u\in N_t\) such that \(u,-u\) are regular normal values
for \(G\) on the interior. This excludes a null set by Sard's theorem.
The reference height, extended over the whole coordinate space,
completes the square as \(h_0=c+bq(z-z_0)\), where
\(b=Rt u_5>0\) and \(z_0\in\Omega^\circ\) is its unique
critical point. Let \(h\) agree with \(\langle G,u\rangle\)
in \(\Omega\) and with \(h_0\) outside. Equality of the full
collar germ makes this a smooth extension. Choose \(L>0\)
larger than both \(|h-c|\) and \(|h_0-c|\) on \(\Omega\).

Write \(z-z_0=(\xi,\eta)\) and \(S=|\xi|^2+|\eta|^2\).
Choose \(S_0\) with \(\Omega\subset\{S<S_0\}\) and \(bS_0/2>L\).
Add the same radial function \(\Psi(S)\) to both heights,
with \(\Psi=0\) up to \(S_0\) and \(\Psi'\) smooth and strictly
increasing from zero to infinity for \(S>S_0\).
The resulting functions are proper and have identical end
sublevel sets at \(c-L,c+L\). A new exterior critical point can
only satisfy \(\eta=0\) and \(2\Psi'(S)=b\). Its critical
value is below \(c-L\), since \(-bS/2+\Psi(S)\) decreases
from its value below \(-L\) at \(S_0\) to that radius.
The compact band between the chosen regular levels therefore
contains only the filling's nondegenerate interior critical points.

The common pair of end sublevels can be computed with
\(h_0+\Psi\). Contracting \(\eta\) to zero decreases that
height and is a deformation of pairs. The function on the negative
plane decreases radially to a minimum below \(-L\) and
then increases without bound. The upper retract is a disk and
the lower one a radial annulus, so the pair has relative
\(H_2\cong\mathbb Z\). Regular-band deformation and single-point
Morse attachment imply that this rank is bounded by the number
of index-two points in the compact band
\cite[Theorems 3.1--3.2 and Remark~3.3]{Milnor1963MorseTheory}.
One may separate coinciding critical values by small bumps
constant near the critical points, preserving their indices.
The auxiliary critical locus below \(c-L\) plays no role.
This proves the count, and reversing the height preserves index two.

Integrate the count over \(N_t\). Its antipode is disjoint
from it. The area formula, counting each sign in its own
normal image, gives \(I_2(G;\Omega)\ge s_t\).
This middle mass is positive, so the strict local comparison
in Lemma~\ref{lem:area} gives the energy lower bound.
\end{proof}

The underlying relative invariant is the degree-two homology of
a pair of fixed end sublevels. A boundary gradient degree alone
would count all Morse indices with alternating signs and would
also allow a definite critical point. More generally a supplied
common exterior extension of two heights with compact regular
end levels fixes a lower bound on the number of interior index-two
points by the rank of its relative \(H_2\) minus the number
known to lie in the auxiliary exterior band. Applying this for a
measurable family of actual ambient directions is what produces
a geometric lower bound.

In the flat limit \(s_t=|\Omega|t^4+O(t^6)\), so
\eqref{eq:scale} and \eqref{eq:normal-area} give the same
leading coefficient as the fixed-germ lower bound. This limit
varies the normalized collar. At any fixed \(t>0\), even
allowing arbitrary immersed interiors leaves the positive
lower bound in Theorem~\ref{thm:fixed}.

\section{An exterior recognition criterion}\label{sec:exterior}

\begin{theorem}[Directional coverage by the reference patches]
\label{thm:exterior}
Let \(M\) be a closed oriented smooth homotopy four-sphere with a
smooth immersion. Let \(P_1,\ldots,P_N\) be pairwise disjoint
closed marked patches modeled on \(\Omega\). Assume that the immersion
equals on a collar of each whole seam a rigidly placed
\(F_{R_j,t_j}\), where \(R_j,t_j>0\).
Let \(C=M\setminus\bigcup_j P_j^\circ\) be the retained exterior
with its induced metric. Let \(N_j\) be the reference one-sign normal
image in actual ambient coordinates, \(s_j=\sigma_4(N_j)\), and
\[
 k(u)=\sum_j{\bf1}_{N_j\cup(-N_j)}(u),\qquad
 D_C=\tfrac12\int_{S^4}(2-k(u))_+\,d\sigma_4(u).
\]
If \(M\) is nonstandard then \(E_+(C)>\tfrac83D_C\).
In particular either inequality
\[
 E_+(C)\le\tfrac83D_C,\qquad
 E_+(C)+\tfrac83\sum_j s_j\le\tfrac83\omega_4
\]
suffices for smooth standardness. These conclusions use
the actual seams and the given exterior metric.
\end{theorem}

\begin{proof}
Replace the interiors by the reference graphs, leaving the
whole exterior fixed. This gives a smooth immersion on the
same \(M\) by equality of maps on collars. In almost every
direction its actual index-two count is \(m_{2,C}(u)+k(u)\).
The formula counts on the domain, regardless of intersections
between immersed pieces. Lemma~\ref{lem:top} forces this sum
to be at least two on a nonstandard \(M\). Integration by
Lemma~\ref{lem:area} shows \(I_2(C)\ge D_C\). Equivalently,
\begin{equation}\label{eq:coverage}
 D_C=\omega_4-\sum_j s_j+
                 \tfrac12\int_{S^4}(k(u)-2)_+\,d\sigma_4(u).
\end{equation}
Each nondegenerate reference seam has middle signature.
Smoothness therefore gives an open middle-signature region on
the exterior side whenever there is a patch. If the family is
empty, the nonstandard count itself gives positive middle
mass. In either case Lemma~\ref{lem:area} gives the strict
exterior bound. The last integral in \eqref{eq:coverage} is
nonnegative, which yields the coarser criterion.
\end{proof}

For a common vertical axis and common horizontal marking,
the upper gnomonic images \(-t_jD\Omega^\circ\) are
nested by star-shapedness. On the double caps
\(\min(2,k)\) is the sum of the two largest indicators.
If \(s_{(1)},s_{(2)}\) are the two largest areas, with zero
for an absent patch, then \(D_C=\omega_4-s_{(1)}-s_{(2)}\).
For separately rotated references, the integral in
\eqref{eq:coverage} retains their actual overlap. For
a finite collection of scaled matching graphs in
Section~\ref{sec:graph}, without the exact fixed reference germs
of this theorem, the simpler consequence is
\(I_2(C)\ge\omega_4-O(\sum_j t_j^4)\) on a nonstandard
closed completion, by the integrated height count and the
patch curvature bounds. This makes the relative-to-global
obstruction explicit when the graph amplitudes tend to zero.

The exterior and seam assumptions are part of the construction
problem. The proofs apply to one immersion on the closed
domain and keep its common exterior fixed. If matching requires
altering that exterior, its new energy enters the comparison.
For positive Yamabe entry, one needs a positive Yamabe constant
for the assembled class as well as the full \(E_+<8\pi^2\)
bound. The exact quadratic patch itself has negative scalar
curvature.

\begin{proposition}[Negative Gauss mass of the exterior]\label{prop:relative-gauss}
Use the hypotheses of Theorem~\ref{thm:exterior}. Choose each
\(N_j\) using the orientation-compatible normal on the graph
and write \(\ell(u)=\sum_j{\bf1}_{N_j}(u)\) and \(S=\sum_j s_j\).
For almost every \(u\), the signed Gauss count on \(C^\circ\)
over \(u\) is \(1-\ell(u)\). With \(m_i(u)\) the height counts after
insertion of the references, the exterior satisfies
\begin{equation}\label{eq:exterior-gauss}
 J(C)=S+I_2(C)+\tfrac12\int_{S^4}
                 (m_0(u)+m_4(u)-2)\,d\sigma_4(u).
\end{equation}
If \(M\) is nonstandard, then
\begin{equation}\label{eq:exterior-shape}
 \begin{aligned}
 J(C)&\ge S+D_C=\omega_4+\tfrac12\int_{S^4}(k(u)-2)_+\,d\sigma_4(u),\\
 B^2(C)&>\tfrac{64}{3}D_C+\tfrac{16}{3}S .
 \end{aligned}
\end{equation}
\end{proposition}

\begin{proof}
The orientation-compatible normal identifies \(TM\) with
\(\nu^*TS^4\) as oriented bundles, so \(2\deg\nu=\chi(M)=2\).
Each reference contributes one positive preimage over \(N_j\)
and zero elsewhere away from the seam image. Subtracting these
contributions gives the signed count on \(C^\circ\).

For almost every height direction, the seams contain no critical
points and the patches contribute only their \(k(u)=\ell(u)+\ell(-u)\)
middle points. The Morse Euler sum therefore gives
\[
 m_1(u)+m_3(u)=m_{2,C}(u)+k(u)+m_0(u)+m_4(u)-2 .
\]
Odd index is equivalent to negative shape determinant at either
normal sign. The area formula identifies the half-integral of
the odd count with \(J(C)\), and the half-integral of \(m_{2,C}\)
with \(I_2(C)\) \cite[\S3, equation (5)]{OntiVlachos2017Hypersurfaces}.
The half-integral of \(k\) is \(S\), giving \eqref{eq:exterior-gauss}.
Every height has a minimum and a maximum, so \(J(C)\ge S+I_2(C)\).
If \(M\) is nonstandard, Theorem~\ref{thm:exterior} and its proof
give \(I_2(C)\ge D_C\) and \(E_+(C)>\frac83D_C\).
The pointwise shape deficit in the immersion paper
\cite[Proposition 2.1]{Baker2026ImmersionEnergy} gives
\(B^2(C)\ge 6E_+(C)+\frac{16}{3}J(C)\), proving
\eqref{eq:exterior-shape}.
\end{proof}

In particular, \(J(C)\ge S\) holds on every smooth type under
the seam hypotheses. If three patches have the same one-sign
image of area \(s\), then \(D_C=\omega_4-2s\). A nonstandard
type requires \(J(C)\ge\omega_4+s\), and almost every direction
in the common one-sign image has at least two negative Gauss
preimages in \(C^\circ\) on any type.

\begin{example}[An exact quadratic exterior on the standard sphere]
\label{ex:exact-exterior}
Translate the unit sphere so its north pole is the origin, and
write its upper chart as the graph of \(h_0(X)=\sqrt{1-|X|^2}-1\).
Fix \(a_j=5je_1\), \(j=1,2,3\), and a smooth cutoff \(\chi\)
supported in \(B_2(0)\) and equal to one on \(B_{3/2}(0)\).
For all sufficiently small \(r>0\), the new chart height
\[
 h_r(X)=h_0(X)+\sum_{j=1}^3\chi(X/r-a_j)
        \bigl(r^2q(X/r-a_j)-h_0(X)\bigr)
\]
has disjoint modification supports compactly contained in the
upper chart. The horizontal differential has rank four, the
change is \(O(r^2)\), and the modified upper sheet stays separate
from the lower one. Thus this is a smooth embedding of \(S^4\).

Let \(P_{j,r}\) be the chart image of \(X=r(a_j+z)\), \(z\in\Omega\).
Since \(\Omega\subset B_{3/2}(0)\), the graph is the translated
\(F_{r,r}\) on each patch and on a neighborhood of its whole seam.
The exterior \(C_r\) is the complement of the interiors of these
three marked chart disks. On all three patches the actual height
levels \(\pm r^2\lambda\) carry the product solid tori and the
side coordinates of Proposition~\ref{prop:trace}, including the
incoming and belt framings. For each \(r\), Sard's theorem permits
choosing \(\delta_0<\lambda<1/2\) so both are regular levels on
the whole sphere.

On the chart supports, \(|Dh_r|=O(r)\) and \(|D^2h_r|=O(1)\).
In fact the bracket and its first two derivatives are \(O(r^2)\),
\(O(r)\), and \(O(1)\), respectively, by the zero constant and
linear jets of \(h_0\) at the origin. The cutoff derivatives of
orders one and two cost \(O(r^{-1})\) and \(O(r^{-2})\). The
resulting shape is uniformly bounded on a set of volume \(O(r^4)\).
The remainder is round and contributes zero to \(E_\pm,B^2,J\).
Consequently \(E_\pm(S^4),B^2(S^4),J(C_r)=O(r^4)\).
The three reference Gauss images are identical by translation,
with area \(s_r=|\Omega|r^4+O(r^6)\) by \eqref{eq:normal-area}.
Thus \(J(C_r)\ge3s_r\) by \eqref{eq:exterior-gauss}, and the
quartic order of the required seam contribution is attained.
The actual exterior satisfies \(E_+(C_r)\le\frac83D_{C_r}\),
where \(D_{C_r}=\omega_4-2s_r\), for small \(r\).

The full induced metric has \(E_+<8\pi^2\) and positive Yamabe
constant for small \(r\). Indeed \(g_r\) is uniformly equivalent
to the round metric, and
\(\|\Scal_{g_r}-12\|_{L^2(g_r)}=O(r^2)\) by the shape and volume
estimates. Compact Sobolev and Hölder give
\[
 \int (6|df|^2+\Scal_{g_r} f^2)\,dV_{g_r}
 \ge (6-Cr^2)\bigl(\|df\|_{L^2(g_r)}^2+\|f\|_{L^2(g_r)}^2\bigr)
 \ge c\|f\|_{L^4(g_r)}^2
\]
with \(c>0\) uniform for small \(r\)
\cite[(1.3)--(1.4) and Theorem 2.3]{LeeParker1987Yamabe}.
\end{example}

The same formula can insert a shrinking cluster in a tangent
graph chart of any fixed smooth immersion. The exterior energy
then converges to that immersion's original full energy, since
both replaced regions have bounded curvature and volume \(O(r^4)\).
Meanwhile \(D_C\to\omega_4\). On a hypothetical nonstandard
type the original full energy lies strictly above
\(\frac83\omega_4\) by the closed theorem, so this procedure
keeps the exterior above the recognition threshold for all
sufficiently small \(r\).

\section{Two comparisons}

\subsection{A spherical two-face transition}

A different two-face model has rotating spherical rather than
solid-torus sections. Set \(F(s,\theta)=(r(s)\theta,z(s))\),
\(\theta\in S^3\), with \(s\) arclength and \(r'=\cos\phi\),
\(z'=\sin\phi\). The principal curvatures in a compatible
orientation are \(p=\phi'\) and \(q_0=\sin\phi/r\) with
multiplicity three. The three opposite-pair sums agree, so
both chiral Weyl blocks vanish. The choice \(p=-q_0/3\)
gives \(r^{1/3}\sin\phi=r_0^{1/3}\) for a waist \(r_0>0\).
The identity \(ds/d\phi=-3r_0\sin^{-4}\phi\) gives infinite
arclength at both angular endpoints and \(r\ge r_0\).
The complete bridge belongs to the rotational
\(H_3=0\) family
\cite[\S1.2.1, equation (7) and Theorem~2(ii)]{Palmas1999RotationHypersurfaces}.
The determinant integral has density
\(\phi'\sin^3\phi\,ds\,d\sigma_3\), and the angle decreases
from \(\pi\) to zero. Thus the signed mass is \(-\omega_4\)
and \(J=\omega_4\), despite the vanishing chiral cost.

On either arm one may replace \(a(r)=(r_0/r)^{1/3}\)
by \(a(r)\chi(r/R)\) on \(R\le r\le2R\), with a nonincreasing
cutoff equal to one and zero near the two ends.
Integrating the outward height derivative
\(a\chi/\sqrt{1-(a\chi)^2}\) gives exact flat end germs
on coaxial parallel planes. The signed determinant is unchanged
by the smoothing, because the two angular endpoints are unchanged.
Monotonicity makes the determinant nonpositive and preserves \(J\).
On the two bands \(|p|+|q_0|=O((r_0/R)^{1/3}/R)\) and their
volume is \(O(R^4)\), so their shape costs are
\(O((r_0/R)^{4/3})\). The spherical model isolates a
negative Gauss charge for one marked transition. It supplies
spherical planar collars with opposite normal ends, whereas the
quadratic construction uses solid tori and positive middle mass.

\subsection{Two actual births on the standard sphere}

The following construction tests the height marking on a single
closed embedding. In local variables \(w=(x_1,x_2,y)\) and \(v\)
put \(Q=-|x|^2+y^2\), choose a cutoff \(\chi\) equal to one
near zero, and set
\[
 P_\tau=Q+v^3+3a^2v-6\tau a^2v\chi(w,v),\quad
 \psi(w,v)=(w,P_0(w,v)),\quad U_\tau=P_\tau\circ\psi^{-1}.
\]
The positive \(v\)-derivative of \(P_0\) makes \(\psi\) a
global diffeomorphism. The function \(U_0\) is the affine
height coordinate and \(U_\tau\) agrees with it off a compact
set. On the fixed transition the differential of \(Q+v^3\)
is bounded away from zero, so sufficiently small fixed \(a>0\)
excludes extra critical points there. For \(\tau=1\) the
inner polynomial has the index-two point \(v=a\) at value
\(-2a^3\) and the index-three point \(v=-a\) at value \(2a^3\).
This is the compactly supported Cerf insertion
\cite[\S3.5, equation (3.5.1)]{BorodzikPowellTeichner2025LinkConcordance}.

Use small disjoint angular charts \(\Theta_j(w)\) on the unit
horizontal \(S^3\), and write the upper unit round sphere over
\(z\in\mathbb R^4\) as \(h_0(z)=\sqrt{1-|z|^2}\).
For \(r\) small choose a common regular latitude
\(h_*=\sqrt{1-R_0^2r^2}\) and normalize the height by
\(T=(h_0-h_*)/r^2\). The two charts have horizontal maps
\[
 z=r\sqrt{R_0^2-2h_*T-r^2T^2}\,\Theta_j(w).
\]
Choose \(R_0\) fixed large enough that the radicand stays
uniformly positive on fixed boxes containing the full supports
and their height ranges. Insert \(h_*+r^2U_\tau(w,T)\) on
each chart with fixed nearby unequal parameters
\(a_1,a_2\), so both index-two levels precede both index-three
levels and all critical values are distinct at the endpoint.
The graphs equal the round map on collars. The horizontal
projection has full rank and remains injective in the upper
graph, and the modified heights stay above the lower hemisphere.
This gives one smooth embedded sphere, with two new index-two
points and two index-three points for the actual vertical height.

In each box the normalized coordinates have bounded first and
second derivatives in \(z/r\), since the boxes stay a
fixed normalized distance from the pole. Thus \(|D_z h|=O(r)\)
and \(|D_z^2h|=O(1)\), uniformly in \(\tau\), and the modified
volume is \(O(r^4)\). The unchanged exterior is round and the
shape stays bounded in the boxes, giving \(E_\pm=O(r^4)\)
along the entire path. This same path gives positive Yamabe
sign for small \(r\). It starts at the round class, and if
an intermediate conformal Laplacian first eigenvalue vanished,
a scalar-flat conformal metric would give
\(E_+=8\pi^2+\frac14\int|\operatorname{Ric}|^2dV\)
by signature and Gauss--Bonnet. This contradicts the path
bound. The identification of the first eigenvalue sign with the
Yamabe sign uses conformal covariance
\cite[\S2]{LeeParker1987Yamabe}.

For clarity, the actual framed link can be read in the cubic
coordinates. Put \(C_a(v)=v^3-3a^2v\) and
\(L_0=3\max_j a_j^3\), with \(a_2/a_1=11/10\).
In the common normalized incoming height \(-L_0\) the \(j\)-th
attaching circle is \(v=a,\ y=0,\ |x|^2=L_0-2a^3\).
A spanning disk in that level is
\[
 y=0,\qquad v_0\le v\le a,\qquad |x|^2=C_a(v)+L_0,
\]
where \(v_0\in(-3a,-2a)\) is its simple left root.
The displayed pieces lie inside the cubic region as the fixed
\(a_j\) are chosen sufficiently small. At the circle the two
normal framing vectors are positive rescalings of \(\partial_v\)
and \(\partial_y\), given by the Morse variables
\((v-a)\sqrt{v+2a},y\). The vector \(\partial_y\) extends
over the disk, normal to it within the level. This gives the
zero framing. Below the first middle critical value the global
level is \(S^3\) by regular transport from an unmodified latitude.
The disks lie in disjoint angular boxes and have disjoint
three-ball neighborhoods in this level. Hence the two circles
form the zero-framed split unlink.

In the intermediate height zero the index-two belt circle
is \(x=0,\ y^2=-C_a(v)\) for \(0\le v\le\sqrt3a\), and
the index-three attaching sphere is \(y=0,\ |x|^2=C_a(v)\)
for \(-\sqrt3a\le v\le0\). The cubic gradient follows the
appropriate stable or unstable directions on these pieces.
They meet transversely once in each box at the origin, because
\((x,y)\) are coordinates in the level there by
\(C_a'(0)\ne0\). There are no cross intersections. Farther
out the incoming and outgoing end levels are exactly round
latitudes. The original angular coordinate and \(T\) give
their product identification and preserve the unchanged
exterior. This realizes an actual subcritical two-middle-handle
height on the known standard type. It leaves the availability
of corresponding exterior data on other types open.

\section{A spectral obstruction on a marked punctured double}
\label{sec:marked-double-exterior}

Let \(W\) be a smooth compact oriented contractible four-manifold
with connected boundary \(Y\). An orientation-preserving diffeomorphism
\(f:Y\to Y\) gives the oriented double \(W\cup_f(-W)\).
Van Kampen and Mayer--Vietoris give a smooth homotopy four-sphere.
Insert \([0,L]\times Y\) using the identity and \(f\) on its two
faces. Collared reparametrization preserves the given smooth type
and is fixed off the seam collars. Mark an interior coordinate ball
on each side, disjoint from these collars. Smooth metrics on the
punctured halves can have arbitrary prescribed product seam
germs and exact unit round product germs at their marked
punctures. Use the pulled-back seam metrics in the corresponding
face markings and interpolate to interior metrics in the complements
of the prescribed germs. The round ends admit conformally flat
caps in their original ball parametrizations.

On a closed four-manifold write \(\Lambda\) for the maximum of
the largest signed eigenvalues of the chiral Weyl blocks. The
ordinary and modified quadratic forms have integrand
\(6|dv|^2+Vv^2\), with \(V=R\) and \(V=R-6\Lambda\),
respectively. On the retained exterior the contribution from a
unit round end is
\(6\langle\tau,\sqrt{1-\Delta_{S^3}}\tau\rangle\) for its trace
\(\tau\), by minimizing on the added conformally flat half
cylinder. This follows by decomposing into eigenfunctions of
\(-\Delta_{S^3}\); each mode minimizes the one-dimensional
quadratic action with prescribed value and decay at infinity.
For the identity double with reflected metrics and corresponding
marked punctures, use all \(H^1\) tests on one punctured half
and leave the trace on the doubled seam free. Positivity of
the corresponding one-half form, including its round-end term,
is equivalent to positivity of the full exterior form.
Restriction gives the sufficient implication by addition,
and even extension gives necessity since it doubles the value.
Orientation reversal exchanges chiralities and preserves
\(\Lambda\). It also makes the full positive chiral energy
the sum of both half energies, with no Weyl energy in the caps.

\begin{proposition}[Conformally product seam obstruction]
\label{prop:marked-double-neck}
Suppose the retained metric on the whole strip \([0,L]\times Y\)
is conformal to \(dt^2+\gamma\), where \(\gamma\) is a smooth
metric on the closed three-manifold \(Y\). Let \(r_{\min}\) be
the least Ricci eigenvalue of \(\gamma\) and use the transverse
first eigenvalues
\[
 \ell_{\rm mod}=\lambda_1(-6\Delta_\gamma+3r_{\min}),\qquad
 \ell_{\rm ord}=\lambda_1(-6\Delta_\gamma+R_\gamma).
\]
If either \(\ell<0\) and \(L\sqrt{-\ell}\ge\pi\sqrt6\), the
corresponding exterior form fails strict positivity, regardless
of the half metrics, puncture caps or twist. In particular, the
ordinary test excludes positive Yamabe sign on every closed
completion that retains this strip. The strip contributes the
exact energy \(\frac L4\int_Y|\operatorname{Ric}_{0,\gamma}|^2
\,dV_\gamma\) to either chirality.
\end{proposition}

\begin{proof}
In an orthonormal product frame whose last three vectors
diagonalize \(\operatorname{Ric}_\gamma\), the only sectional
curvature in the opposite pair \(0i,jk\) is
\(K_{jk}=R_\gamma/2-r_i\). Project onto either chiral basis and
subtract the scalar contribution \(R_\gamma/12\).
The Weyl eigenvalues are \(R_\gamma/6-r_i/2\), giving the
modified potential and energy identity. The lowest Dirichlet
value of the product strip is \(\ell+6\pi^2/L^2\), by the
Poincaré inequality in \(t\) and the transverse Rayleigh
inequality. It is attained by \(\sin(\pi t/L)\phi(y)\) for a
first transverse eigenfunction \(\phi\). Its zero extension
is a nonzero \(H^1\) test on the retained exterior, with zero
trace on every marked round end. The scalar quadratic form
transfers under multiplication of the test by the inverse
conformal factor, since the test has zero trace on both strip
faces. The added potential transfers because \(\Lambda\) has
conformal weight \(-2\). The Weyl energy of the strip is
conformally invariant in four dimensions. A warped collar
\(ds^2+a(s)^2\gamma\) is included with \(L=\int ds/a(s)\).
\end{proof}

The same product calculation gives a shorter bound for fixed
half metrics. Normalize \(\phi\) in \(L^2(\gamma)\), take one
test on each half with that trace in the common cylinder
markings and zero trace at the marked punctures, and let \(H\)
be the sum of their form values off the inserted strip. Put
\(k=\sqrt{-\ell/6}\). When \(0<L<\pi/k\), the intervening test
\[
 \frac{\cos(k(t-L/2))}{\cos(kL/2)}\phi(y)
\]
has both required traces. Integration by parts gives strip
value \(-12k\tan(kL/2)\), so \(H\le12k\tan(kL/2)\) suffices
for failure. The bound includes the cost of realizing the
two traces in the fixed halves. In the second copy the
actual trace in its own marking is \(\phi\circ f^{-1}\).

\begin{example}[The Mazur seam]\label{ex:mazur-spectral-seam}
For \(n=1\), the lower half of
\href{https://arxiv.org/pdf/2403.17356v2#page=35}{Isoshima's
Figure 37} is the Mazur manifold \(W^-(0,3)\), with the cut
after its first two-handle. The identity double is standard
by Isoshima's Theorem 3.3. The punctures can be marked in
the two corresponding zero-handles. The actual middle
cut is \(Y=\Sigma(2,5,7)\), using
\href{https://math.berkeley.edu/~kirby/papers/Akbulut%20and%20Kirby%20-%20Mazur%20manifolds%20-%20MR0544597.pdf}{Akbulut--Kirby,
Proposition 1 and Theorem 2(1)}.

For every \(\gamma\) on this \(Y\), both transverse
eigenvalues in Proposition~\ref{prop:marked-double-neck}
are negative. Indeed, \(1/2+1/5+1/7<1\), so
\href{https://webhomes.maths.ed.ac.uk/~v1ranick/papers/milnbries.pdf}{Milnor,
Theorem 2.7 and Corollary 6.6}, gives a finite oriented
circle-bundle cover \(Z\) of \(Y\) over a closed orientable
surface of genus at least two. The preimage of an essential
simple closed curve is a torus with injective fundamental
group, by the homotopy sequence of the fibration. The
projection onto the surface group also gives \(b_1(Z)\ge4\).
Any nonnegative scalar curvature metric on \(Y\) would lift
to \(Z\) and force that metric to be flat by
\href{https://archive.ymsc.tsinghua.edu.cn/pacm_download/59/11113-Shing-Tung_Yau_216.pdf}{Schoen--Yau,
Theorem 3, p.~2567}. Parallel harmonic one-forms would give
\(b_1(Z)\le3\), a contradiction. Thus the lowest eigenvalue
of \(-8\Delta_\gamma+R_\gamma\) is negative, since a
nonnegative value would give a nonnegative scalar curvature
metric by conformal change with the positive first
eigenfunction \cite[eq.~(1.1)]{LeeParker1987Yamabe}. Both
transverse forms are bounded above on that test by this
three-dimensional conformal form.

For an explicit collar take the metric descending from the
left-invariant metric on
\(\widetilde{\mathrm{SL}}(2,\mathbb R)\) with cyclic
orthonormal bracket constants \((a,a,-a)\), \(a>0\).
Use Milnor's Brieskorn quotient model (Corollary 6.5)
and his
\href{https://core.ac.uk/download/pdf/82428733.pdf}{Ricci
formula (Theorem 4.3)}. The eigenvalues of Ricci are
\((-3a^2/2,-3a^2/2,a^2/2)\). The modified and ordinary
transverse potentials are constant, with the corresponding
obstructions at \(aL\ge 2\pi/\sqrt3\) and
\(aL\ge\pi\sqrt{12/5}\). Either chiral charge is
\((2/3)a^4L\operatorname{Vol}_{\gamma}(Y)\).
This supplies a restriction on the conformally product
construction at the specified marked cut. Shorter collars
require an estimate on the half metrics or further geometric
data.
\end{example}

\section{Short marked product seams and the cost of the halves}
\label{sec:short-marked-seam}

Keep the retained punctured halves \(X_0,X_1\) and their metrics
fixed. Use exact unit round collars at their marked punctures and
fix the seam face markings \(j_i:(Y,\gamma)\to\partial_{\rm seam}X_i\).
The seam germs are exact products. The inserted product has length
\(L>0\) and is glued by \(j_0,j_1\). In the double notation these
maps are the identity and the given orientation-preserving \(f\),
so the trace at the second face in its original marking is
\(\psi_1\circ f^{-1}\). The second half's orientation is reversed
in the joined four-manifold. Collar insertion and capping the
round punctures in their original coordinates give the same
closed smooth type as the length-zero gluing.

For either of the scalar potentials \(V\) from
Section~\ref{sec:marked-double-exterior}, define the form on a
retained half by
\[
 q_i(u)=\int_{X_i}(6|du|^2+Vu^2)\,dV
 +6\sum_b\langle\tau_bu,\sqrt{1-\Delta_{S^3}}\,\tau_bu\rangle,
 \qquad
 d_i[\psi_i]=\inf_{u|_{\partial_{\rm seam}X_i}
                         =\psi_i\circ j_i^{-1}}q_i(u).
\]
For this definition of the trace minimum require that \(q_i\)
be positive on every nonzero test of zero seam trace. Failure
on that subspace gives a global nonpositive test by zero
extension. The compact \(L^2\) embedding and bounded potential
give coercivity on the subspace when it is positive.
The minimizer for a fixed trace is the unique weak solution
of \((-6\Delta+V)u=0\) with prescribed seam trace and
the natural round-end condition
\(6\partial_\nu u+6\sqrt{1-\Delta_{S^3}}\tau_bu=0\) on each
puncture face. Thus \(d_i\) is the pairing of the prescribed
trace with the outward conormal derivative \(6\partial_\nu u\),
in the indicated markings.

\begin{proposition}[Exact trace test]\label{prop:short-seam-trace}
Let \(\ell_1\) be the first eigenvalue of the transverse operator
\(A=-6\Delta_\gamma+V_\gamma\) from
Proposition~\ref{prop:marked-double-neck}. Suppose
\(\ell_1+6\pi^2/L^2>0\). In an \(L^2(\gamma)\)-orthonormal
eigenbasis of \(A\), denote the coefficients of the seam traces
\(\psi_i\) by \(x_{ij}\). For an eigenvalue \(\ell\), put
\(k=\sqrt{|\ell|/6}\). The modal matrices of the exact cylinder
trace form \(c_L\) are
\[
 \begin{aligned}
 C_L(\ell>0)&=6k
 \begin{pmatrix}\coth(kL)&-\operatorname{csch}(kL)\\
                 -\operatorname{csch}(kL)&\coth(kL)\end{pmatrix},\\
 C_L(0)&=\tfrac6L\begin{pmatrix}1&-1\\-1&1\end{pmatrix},\\
 C_L(\ell<0)&=6k
 \begin{pmatrix}\cot(kL)&-\csc(kL)\\-\csc(kL)&\cot(kL)\end{pmatrix}.
 \end{aligned}
\]
The capped form is strictly positive if and only if both
zero-seam-trace half forms are strictly positive and
\[
 d_0[\psi_0]+d_1[\psi_1]
 +c_L[\psi_0,\psi_1]>0,\qquad
 c_L=\sum_j (x_{0j},x_{1j}) C_L(\ell_j)(x_{0j},x_{1j})^{\mathsf T}
\]
for every nonzero \(H^{1/2}(Y)\) trace pair. For the ordinary
form this is equivalent to positive Yamabe sign.
\end{proposition}

\begin{proof}
The cylinder minimizer for each trace pair solves
\(6b''=\ell b\) on the mode of eigenvalue \(\ell\). Its action
is \(6[b b']_0^L\), which gives the matrices by the specified
endpoint values. The zero-endpoint cylinder form is coercive
by the length hypothesis and the two Rayleigh inequalities.
Subtracting the three fixed-trace minimizers gives an exact
decomposition into the trace form and the three zero-trace
quadratic forms. The closed form and the round-end reduction
have the same strict sign. On the closed compact manifold a
strictly positive form with bounded potential has an attained
positive first \(L^2\) eigenvalue. This and the Sobolev inequality
give the positive normalized quotient.
\end{proof}

For the reflected identity double use the free trace at the
midpoint. The one-half trace test becomes \(d_h+n_L\), where
the eigenvalues of \(n_L\) are \(6k\tanh(kL/2)\), \(0\), and
\(-6k\tan(kL/2)\) on positive, zero, and negative transverse
eigenvalues, respectively. This follows by imposing
\(b'(L/2)=0\) on the half-cylinder minimizer. Even extension and
restriction give equivalence with the identity double test.
If the same retained half is used on both sides and \(f\) is
a \(\gamma\)-isometry, cutting the twisted double at the
midpoint gives two copies of the same free-trace problem.
The identity sign then implies the twisted sign.
For example, let \(q_h\ge c\|\cdot\|_{H^1}^2\) on the unextended
half and let \(C_{\rm tr}\) be its \(H^1\)-to-\(L^2\) seam
trace bound. If \(\ell_1<0\), both forms pass for
\(L<2k_0^{-1}\arctan(c/(6k_0 C_{\rm tr}^2))\) for the selected
potential, where \(k_0=\sqrt{-\ell_1/6}\). This is because the
negative part of \(n_L\) is bounded above by
\(6k_0\tan(k_0L/2)\) in \(L^2\).

In the original orientations of the two retained metrics
the exact closed energy is
\begin{equation}
 \label{eq:short-seam-energy}
 E_\pm(L)=E_\pm(h_0;X_0)+E_\mp(h_1;X_1)
       +\frac L4\int_Y|\operatorname{Ric}^0_\gamma|^2\,dV_\gamma.
\end{equation}
The round caps have zero Weyl curvature and the half terms
include any retained interpolation. With the same metric on
both halves, the two chiral energies and their dependence on
length agree for identity and compatible twisted gluing.

\begin{proposition}[Signature cost of the homogeneous Mazur germ]
\label{prop:short-seam-aps}
On the homogeneous metrics \(\gamma_a\) of the marked
\(\Sigma(2,5,7)\) in Example~\ref{ex:mazur-spectral-seam},
let the two retained product face germs be pulled back from
\(\gamma_a\) by the actual face markings. If the capped halves
are contractible, either closed chirality has energy at least
\[
 12\pi^2|\eta_{\rm sig}(Y,\gamma_a)|
       +\frac{484\pi^2}{105}\,aL .
\]
Here \(\eta_{\rm sig}\) is the unreduced eta value of the
odd signature operator on even forms. For this metric
\(|\eta_{\rm sig}|=65/84\), so the energy exceeds
\(8\pi^2\) for every positive \(L\).
\end{proposition}

\begin{proof}
For the product-boundary capped half the signature is zero.
The boundary signature theorem therefore gives
\[
 E_+(h_i;X_i)-E_-(h_i;X_i)
  =12\pi^2\eta_{\rm sig}(\partial X_i,\gamma_a),
\]
where the boundary on the right means the seam face after
capping the punctures. The chiral density is in the
Hilbert--Schmidt operator convention and the eta value has
the outward-normal-first orientation. Capping adds no Weyl
energy. The sum of the two half energies is at least the
absolute value of their difference on each half, and the
closed signature vanishes. These facts and
\eqref{eq:short-seam-energy} give the stated lower bound.
\cite[Theorem (4.14)]{APS1975SpectralI}
\cite[\S2]{BiquardHerzlichRumin2007Diabatic}

We give the metric normalization and eta calculation. The
Seifert direction is the third direction in the cyclic
Milnor frame. The base has curvature \(-a^2\), since
\(K(e_1,e_2)=-7a^2/4\) and the submersion adds \(3a^2/4\).
The oriented triangle quotient has area \(22\pi/(70a^2)\).
Milnor's central extension has generators \(g_i\) with
\(g_i^{\alpha_i}=g_1g_2g_3=c\), where \(c\) is a full
rotation and \(\alpha_i=(2,5,7)\).
Its abelianization is cyclic of order \(11\) generated by \(c\).
Milnor's commutator model thus has regular fiber length
\(22\pi/a\). The volume is \(242\pi^2/(35a^3)\).
The strip charge in the proposition now follows from the
Ricci eigenvalues in Example~\ref{ex:mazur-spectral-seam}.
\href{https://webhomes.maths.ed.ac.uk/~v1ranick/papers/neuwir3.pdf}{Milnor, Theorem 3.1 and Corollaries 3.2 and 6.5}.

Use the unit-period-\(2\pi\) connection
\(\theta=(a/11)e^3\) and orientation \(\theta\wedge d\theta>0\).
The degree is \(d=-\int d\theta/(2\pi)=-1/70\) on the base.
In the abelianization \(g_i\equiv m_i c\) with
\(m_i=(6,9,8)\). The positive base rotation is lifted in the
commutator by \(g_i c^{-m_i}\), whose fiber angle in the
\(2\pi\) normalization is \(2\pi(1/\alpha_i-m_i)/11\).
Hence the local action weights of the orbifold bundle are
\((1-\alpha_i m_i)/11\bmod\alpha_i=(1,1,2)\).
To fix the sign in the adiabatic limit, apply the signature
theorem to the disk orbifold of this Seifert line with the
cone points excised. The oriented inner boundaries are the
negatives of the round lenses \(L(\alpha_i,\beta_i)\).
The rational intersection form has a single negative direction
from the zero section of self-intersection \(d\), and removing
rational homology ball neighborhoods preserves this form.
Use a connection metric with fiber
\(ds^2+F_\epsilon(s)^2\theta^2\), where \(F_\epsilon=s\)
near zero and \(F_\epsilon=\epsilon\) near the outer face.
At fixed \(\epsilon\) the metric can be flattened in shrinking
orbifold normal balls with bounded curvature, and a radial
conformal change then gives round product inner boundaries at
vanishing signature density cost. The outer Pontryagin
integral tends to \(d/3\). Indeed the horizontal bundle and
line reference connections have vanishing Pontryagin four-form.
At the outer boundary their line block has entry \(-\theta\),
the Levi--Civita line block is zero, and the horizontal and
mixed block differences tend to zero. The outer transgression
of \(\operatorname{tr}(\Omega^2)\) tends to
\(2\theta\wedge d\theta\), giving \(d/3\) after multiplication
by \(-1/(24\pi^2)\) and integration.
The APS eta value on the positively oriented round lens is
\(-4s(\beta_i,\alpha_i)\), so the inner boundary orientation
gives
\[
 \eta_{\rm ad}=1+d/3-4\sum_i s(\beta_i,\alpha_i)
              =-19/210.
\]
The three sums are \(0,1/5,1/14\). The sign also follows from
the residue sum in the adiabatic calculation.
\cite[Theorem (4.14)]{APS1975SpectralI}
\cite[Proposition (2.12)]{APS1975SpectralII}
\cite[equations (75)--(77)]{BiquardHerzlichRumin2007Diabatic}

For an independent finite-metric check, vary just the
third coframe to \(v e^3\). The Levi--Civita connection has
\(\omega_{12}=-a(1+v^2/2)e^3\),
\(\omega_{23}=av e^1/2\), and \(\omega_{31}=av e^2/2\).
The Chern--Simons trace
\(\operatorname{tr}(\omega\wedge d\omega+2\omega^3/3)\)
is \(-a^3(2+2v^2+v^4)e^{123}\).
Apply the boundary signature theorem to a product-ended
interpolation cylinder. Its signature is zero and
\(p_1=-\operatorname{tr}(\Omega\wedge\Omega)/(8\pi^2)\).
Thus \(\eta(v)-\eta_{\rm ad}=
 a^3\operatorname{Vol}_{\gamma_a}(Y)(2v^2+v^4)/(24\pi^2)\),
where the limit uses \(v\to0\).
At \(v=1\) the correction is \(121/140\), giving \(65/84\)
in the chosen orientation. Reversing orientation changes the
sign. The obstruction depends on the prescribed homogeneous
boundary metric and survives arbitrary shortening of its
exact product seam.
\end{proof}

The same computation permits an exact change of the marked
cross-section metric. On the same quotient, scale the third
coframe by \(v>0\) and keep the horizontal coframe fixed.
The resulting Ricci eigenvalues are
\((-a^2(1+v^2/2),-a^2(1+v^2/2),a^2v^2/2)\), from the
cyclic brackets \((a/v,a/v,-av)\). Its volume is
\(v\operatorname{Vol}_{\gamma_a}(Y)\). Thus
\[
 \eta_{\rm sig}(v)=\frac{121(2v^2+v^4)-38}{420},\qquad
 E_\pm(\mathrm{strip};v)=\frac{121\pi^2}{105}\,aL\,v(1+v^2)^2 .
\]
At \(v^2=\sqrt{159}/11-1\) the boundary signature cost
vanishes. The two transverse potentials are still negative
constants. For this altered metric the ordinary trace test
and the charged interior energy remain the tests for entry,
with retained halves that meet the new germ.

For a further marked presentation take the Mazur cork and
its nonextendible boundary involution in the diagrams of
Auckly--Kim--Melvin--Ruberman. Their Figure~2 gives smooth
standard spheres for both identity and twisted gluings.
Averaging a boundary metric over the involution gives an
invariant product germ and permits the same chosen half
metric to be used on each side in the identity and twisted
constructions. The sign comparison and
\eqref{eq:short-seam-energy} then apply without extending the
involution. This gives a comparison on two actual smooth
gluings. Realizing the ordinary sign and the strict energy
ceiling on one such fixed half remains a geometric premise.
\cite[\S1 and Figures 1--2]{AucklyKimMelvinRuberman2017Equivariant}



\section{Outstanding challenges}

The first construction problem is to realize the two prescribed
solid-torus markings and common corner in an immersion of
the retained exterior of a given smooth presentation, with a
quantitative energy bound measured on that same exterior.
The fixed-collar estimate identifies the cost of substituting
an interior once this geometry is supplied. To use flattening
instead, one would need a family of exteriors whose normalized
germs vary with the small slopes and whose geometry remains
controlled up to the seams.

For several handles the directional formula uses the actual
rotations and overlap of all reference normal images. A bound
for one patch gives little control of a collection unless its
embeddings and their common exterior are specified together.
A different family of reference germs could be tested by
computing compact sublevel pairs for a set of ambient directions
and determining their exterior critical counts. This would
extend the relative argument to the extent that the required
height and collar data can be obtained.

There are two further global steps. A positive-entry construction
must control the closed conformal Laplacian as well as the total
chiral energy. A height-recognition construction must identify the
actual reduced surgery diagram in one selected subcritical
direction. The standard example supplies a split diagram and
one convenient direction. Neither the index count nor localized
energy controls the isotopy class of a prescribed attaching knot
or the existence of qualifying data on arbitrary smooth
homotopy four-spheres.

\paragraph{Provenance.}
Prepared with OpenAI Codex assistance under the bounded SP4-MASTER-E51a
editorial assignment, with GPT-6.2 Sol/ultra requested.
Actual serving model and reasoning effort are unavailable.
The draft assembles repository results E32, E37, E42, E49 and E50
at commit \texttt{483ad3cef5a0af4328c5d35cabbc1f5224d26b1e}.
The exterior Gauss budget and exact quadratic exterior were added
for SP4-MASTER-E53b, prepared for Edward B. Baker III on
25 November 2026 from checkout \texttt{b5de9758f16e}, with the
actual serving model and reasoning effort unavailable.
The accompanying manuscript README records the evidence and input versions.

The 25 November 2026 revision was prepared for Edward B. Baker III
with OpenAI Codex assistance from the manuscripts at
\texttt{b5ce9f481a76}. The accompanying research assessment records
the E59--E62 derivations, independent readers and source generations.

\bibliographystyle{unsrturl}
\bibliography{references}
\end{document}
